\chapter{Mathematical and Graphical Toolkit}
\label{chap:maths}

\section{Learning objectives}

After completing this chapter, you should be able to:
\begin{itemize}
  \item solve the linear equations used in introductory microeconomics;
  \item interpret functions, graphs, slopes, and intercepts;
  \item calculate and interpret ordinary and partial derivatives; and
  \item use marginal analysis to identify an optimum.
\end{itemize}

\section{Why mathematics matters in microeconomics}

Microeconomics uses mathematics to make an argument precise. A diagram can show that a higher price is associated with a lower quantity demanded; an equation states exactly how the two variables are related. A derivative tells us how rapidly an outcome changes when one variable changes. Mathematics is not an end in itself: it is a language for describing choices, constraints, and trade-offs.

Every result should be read in three ways: verbally, algebraically, and graphically. Ask what the economic question is, what equation represents it, and what the relationship looks like in a diagram.

\section{Variables, parameters, and functions}

A \keyterm{variable} is a quantity that can take different values. In consumer theory, $x$ and $y$ may denote quantities of two goods; in a market, $p$ may denote price and $q$ quantity. A \keyterm{parameter} is treated as fixed within a particular problem, even though it may differ across situations. Income, $I$, and a price, $p_x$, are often parameters when we study a consumer's choices.

A \keyterm{function} describes how one variable depends on one or more others. For example,
\[
q=100-2p
\]
is a demand function. It states that quantity demanded $q$ depends on price $p$. The number $100$ is the intercept and $-2$ is the slope.

\section{Linear equations, slopes, and intercepts}

Many introductory models use linear equations:
\[
y=a+bx,
\]
where $a$ is the vertical intercept and $b$ is the slope. The intercept is the value of $y$ when $x=0$. The slope measures the change in $y$ associated with a one-unit increase in $x$.

Consider inverse demand:
\[
p=50-2q.
\]
When $q=0$, price is $50$, so the vertical intercept is $50$. The slope is $-2$: an additional unit sold is associated with a price two euros lower. Setting $p=0$ gives $q=25$, the horizontal intercept.

To solve an equation, isolate the unknown. If $300-5x=0$, then $5x=300$, and therefore $x=60$. Check an answer by substituting it into the original equation.

The slope between two points is
\[
\text{slope}=\frac{\Delta y}{\Delta x}.
\]
For a straight line, the slope is the same everywhere. A positive slope indicates that $y$ increases when $x$ increases; a negative slope indicates that $y$ decreases when $x$ increases.

\section{Graphs and economic diagrams}

A graph is a visual representation of a relationship. In the usual convention, the horizontal axis measures the independent variable and the vertical axis the dependent variable. Economics sometimes reverses this convention: demand is often drawn with quantity on the horizontal axis and price on the vertical axis because it makes market diagrams easier to read.

When reading a graph, first identify the variables and units on both axes. Then ask whether a movement is caused by a change in a variable shown on an axis, or by a change in another factor that shifts the curve. A price change produces a movement along a demand curve; a change in income, tastes, or the price of a related good may shift the curve.

\section{Derivatives and marginal change}

The \keyterm{derivative} of a function measures its instantaneous rate of change. If $f(x)$ is an economic outcome, the derivative $f'(x)$, or $\frac{df}{dx}$, measures how $f$ changes when $x$ changes by a very small amount.

For a power function,
\[
\frac{d}{dx}x^n=nx^{n-1}.
\]
Thus,
\[
\frac{d}{dx}x^2=2x
\qquad\text{and}\qquad
\frac{d}{dx}\left(\frac{3}{4}x^3+7x^2-5x\right)=\frac{9}{4}x^2+14x-5.
\]

In economics, derivatives are usually called \keyterm{marginal} values. If $C(q)$ is total cost, then $C'(q)$ is marginal cost: the additional cost of producing one more unit, approximately. If $U(x)$ is utility, then $U'(x)$ is marginal utility.

\section{Partial derivatives}

Economic choices often involve more than one variable. A utility function may be $u(x,y)$, where $x$ and $y$ are two goods. A \keyterm{partial derivative} measures the effect of changing one variable while holding the other variables fixed.

For
\[
f(x,y)=2x^2+3xy+y^2-5,
\]
the partial derivative with respect to $x$ is
\[
\frac{\partial f}{\partial x}=4x+3y.
\]
When differentiating with respect to $x$, treat $y$ as a constant. The partial derivative with respect to $y$ is
\[
\frac{\partial f}{\partial y}=3x+2y.
\]

\section{Optimisation}

Many microeconomic problems ask an agent to choose the best feasible option. The structure is always similar:
\begin{enumerate}
  \item identify the objective, such as maximising utility or profit;
  \item identify the constraint, such as a budget or available technology;
  \item express the objective in terms of the choice variable or variables;
  \item find candidate optima using a marginal condition; and
  \item check that the candidate is feasible and is a maximum rather than a minimum.
\end{enumerate}

Consider a consumer with utility $u(x,y)=xy$, income $I=100$, and prices $p_x=10$ and $p_y=5$. The budget constraint is
\[
10x+5y=100,
\]
or $y=20-2x$. Substituting into utility gives
\[
u(x)=x(20-2x)=20x-2x^2.
\]
The first-order condition is
\[
\frac{du}{dx}=20-4x=0,
\]
which gives $x=5$. The budget constraint then gives $y=10$. Since $\frac{d^2u}{dx^2}=-4<0$, this bundle maximises utility.

Correct algebra is necessary but not sufficient. A solution must make economic sense: quantities cannot be negative, choices must satisfy the constraint, and units must be consistent. Always conclude a calculation with an economic interpretation.

\section*{Chapter Review}

\subsection*{Key concepts}
\begin{itemize}
  \item \textbf{Function}: a rule linking a dependent variable to one or more inputs.
  \item \textbf{Slope}: the change in one variable associated with a one-unit change in another.
  \item \textbf{Derivative}: the instantaneous rate of change of a function.
  \item \textbf{Marginal value}: the economic interpretation of a derivative.
  \item \textbf{Partial derivative}: the effect of changing one variable while holding others fixed.
  \item \textbf{Constraint}: a limit on the feasible choices of an economic agent.
  \item \textbf{Optimum}: the best feasible choice according to the objective.
\end{itemize}

\subsection*{Key tools}
\begin{itemize}
  \item Linear equation: $y=a+bx$.
  \item Slope: $\Delta y/\Delta x$.
  \item Power rule: $\frac{d}{dx}x^n=nx^{n-1}$.
  \item An interior optimum is a candidate point where the first derivative equals zero.
  \item A negative second derivative indicates a local maximum.
\end{itemize}

\section*{Exercises}

\subsection*{Conceptual questions}
\begin{enumerate}[label=\arabic*.]
  \item Distinguish a variable from a parameter. Give one example of each from consumer theory.
  \item Explain the economic meaning of the slope of an inverse demand curve.
  \item What does a derivative measure? Why are derivatives often called marginal values in economics?
\end{enumerate}

\subsection*{Analytical exercises}
\begin{enumerate}[label=\arabic*.]
  \item Solve: $300-5x=0$; $\frac{120}{x}=6$; and $(x+1)^2=25$.
  \item Calculate the derivative of: (a) $f(x)=x^2$; (b) $f(x)=\frac{3}{4}x^3+7x^2-5x$; and (c) $f(x)=2x^{35}+3x^{27}-5x^{20}+74x-325$.
  \item Let $f(x,y)=x^2y-xy-6x$. Find $\frac{\partial f}{\partial x}$ and $\frac{\partial f}{\partial y}$.
\end{enumerate}

\subsection*{Graphical exercises}
\begin{enumerate}[resume,label=\arabic*.]
  \item Consider the inverse demand curve $p=60-3q$. Find both intercepts, calculate the slope, and draw the curve with clearly labelled axes.
  \item The budget line is $2x+y=10$. Find the intercepts and draw it. How does the line change if income rises and the new constraint is $2x+y=16$?
  \item Explain, using a diagram, the difference between a movement along a demand curve and a rightward shift of demand.
\end{enumerate}

\subsection*{Optimisation exercises}
\begin{enumerate}[resume,label=\arabic*.]
  \item A consumer has income $I=200$, prices $p_x=20$ and $p_y=10$, and utility $u(x,y)=xy$. Solve the consumer's problem by substituting the budget constraint into the utility function.
  \item Repeat Exercise 7 when utility is $u(x,y)=x^{2/3}y^{1/3}$, income is $100$, $p_x=10$, and $p_y=5$. Interpret the result.
  \item A firm has total revenue $R(q)=50q-2q^2$ and total cost $C(q)=10+2q$. Find profit as a function of output and calculate the profit-maximising quantity.
\end{enumerate}

\subsection*{Exam-style multiple-choice questions}
\begin{enumerate}[resume,label=\arabic*.]
  \item If $f(x)=3x^2-4x+7$, then $f'(x)$ is: (a) $6x-4$; (b) $3x-4$; (c) $6x+7$; or (d) $3x^2-4$.
  \item In $p=100-5q$, the number $-5$ is: (a) the vertical intercept; (b) the slope; (c) the quantity demanded when price is zero; or (d) marginal cost.
  \item A partial derivative with respect to $x$ measures: (a) the effect of changing every variable simultaneously; (b) the effect of changing $x$ while holding other variables fixed; (c) the total value of $x$; or (d) the slope of every possible budget line.
\end{enumerate}

\subsection*{PPE discussion question}

\noindent Economic models simplify reality. When is simplification intellectually useful, and when might it obscure important questions about power, justice, or human behaviour? Illustrate your answer with one microeconomic example.
